\documentclass{article}
\usepackage{amsmath,amssymb}
\usepackage{booktabs}
\usepackage{hyperref}
\usepackage{geometry}
\geometry{margin=1in}

\newcommand{\pos}[1]{\left(#1\right)^{+}}
\setlength{\emergencystretch}{3em}

\title{The Cost of Wordplay: an Educational Programming Platform}
\author{Issy Chen, Siran Luo}
\date{}

\begin{document}
\maketitle

\begin{abstract}
Wordplay is an educational programming platform hosted at the University of Washington with variable cloud charges from Claude LLM calls and Firebase services, both subject to daily free usage quotas. We build a discrete cost model to predict monthly platform expenses given user counts and student activity rates. We use deterministic calculation as a baseline and Monte Carlo simulation to account for random variation in user behavior. Our small test case demonstrates nonlinear cost growth when usage exceeds free tiers, identifies major cost drivers, and lays out a plan to validate predictions against real billing data.
\end{abstract}

\section{Problem Description}
Wordplay’s cloud infrastructure generates charges for LLM API calls, database operations, storage, and data transfer. Every resource has a daily free quota that resets overnight; only excess usage produces fees. The platform team currently lacks a reliable way to forecast how rising student activity will affect monthly spending.

\noindent \textbf{Our goal:} Build a transparent mathematical model to predict Wordplay’s total monthly cloud cost from user count and activity parameters, and pinpoint the dominant cost drivers.

We seek to answer three central questions:
\begin{enumerate}
    \item For a given number of active students and typical usage patterns, what is the expected monthly cloud cost?
    \item How strongly does total cost respond to changes in editing time, project saves, and AI feature usage?
    \item What activity thresholds push the platform past its free quotas and trigger steep cost increases?
\end{enumerate}

Users and resource events are discrete, and student behavior is random, so simple linear approximations fail to capture cost uncertainty. We adopt an incremental modeling approach, starting with simplifying assumptions and adding complexity, with later validation against real billing data.

\section{Simplifications and Assumptions}
The full cloud cost system for Wordplay couples many interacting, partially documented resource charges. To construct a tractable model within the scope of this project, we introduce the following simplifications, each with its associated justification or limitation.
\begin{enumerate}
    \item We use fixed baseline estimates of student behavior rates (writes per editing hour, save frequency, LLM token consumption) provided by the Wordplay team, rather than fitted values from platform logs. These point estimates are reasonable starting values, but they may not reflect real variation across classrooms or student groups.
    \item We compute Firestore costs on a daily basis and apply the daily free quota separately for each day, consistent with cloud provider rules. We do not aggregate usage across the whole month before subtracting free allowance. This simplification correctly captures the daily reset rule, but it requires us to iterate day-by-day in our calculations.
    \item We exclude secondary cost components that are not yet verified or quantified in the Wordplay documentation, such as CPU-second charges and authentication overage fees. This reduces the number of parameters we need to track, yet it means our model will slightly underpredict total cost if those omitted charges become meaningful at higher usage.
    \item In our deterministic baseline, every student’s activity rate is identical. In the stochastic simulation, we sample activity from simple parametric distributions. We do not model correlation between students, such as synchronized class activity during scheduled lessons. This ignores clustered usage spikes that could push daily usage over thresholds more often than independent sampling predicts.
    \item We treat collaborative gallery and shared-project sizes as exogenous input parameters instead of modeling how students naturally form shared projects. This keeps the state space small, but the model cannot endogenously predict how collaboration patterns emerge as the user base grows.
\end{enumerate}

All these simplifications let us build, implement, and test the cost model incrementally, as recommended in the project feedback. Each assumption creates a clear limitation: our outputs are conditional on these choices, and results should be interpreted with those caveats in mind. Later refinements could relax some assumptions if log data or additional billing records become available.


All costs are in US dollars. Prices are for Claude's first-party API and
Google Cloud us-central1. Per-hour activity counts marked ``estimate'' are
assumptions, not measurements.

\paragraph{Notation.} The positive part of $x$ is
\[
\pos{x} = \max(x, 0),
\]
so usage inside a free tier costs nothing.

\section{Claude}

\subsection{Variables}

\begin{center}
\begin{tabular}{ll}
\toprule
Symbol & Meaning \\
\midrule
$I$ & Input tokens in a Haiku call \\
$O$ & Output tokens in a call \\
$U$ & Uncached input tokens in an Opus call \\
$W_c$ & Tokens written to the prompt cache \\
$R_c$ & Tokens read from the prompt cache \\
$T$ & Thinking tokens, billed as output \\
$n$ & Strings being checked \\
$k$ & Chunks of up to 40 strings \\
\bottomrule
\end{tabular}
\end{center}

\subsection{Translation: \texttt{getLLMTranslations} (Haiku 5.5)}

\[
C_{\text{translate}} = \frac{0.10\,I + 0.50\,O}{10^{6}}
\]

\subsection{Reading-level check: \texttt{analyzeLocalization} (Opus 5.5)}

\[
C_{\text{chunk}} =
\frac{4\,U + 5\,W_c + 0.20\,R_c + 20\,(O + T)}{10^{6}}
\]

\[
k = \left\lceil \frac{n}{40} \right\rceil
\]

With the estimates $U = 1500$, $W_c = R_c = 2000$ and $O + T = 5000$:

\[
C_{\text{check}} \approx 0.116 + 0.106\,(k - 1)
\]

\section{Firebase}

\subsection{Variables}

\begin{center}
\begin{tabular}{ll}
\toprule
Symbol & Meaning \\
\midrule
$W$ & Firestore document writes in a day** \\
$R$ & Firestore document reads in a day** \\
$D$ & Firestore document deletes in a day** \\
$G$ & Firestore storage in GiB \\
$N$ & Students in one gallery \\
$H$ & Hours of editing** \\
$S$ & Project saves per student per hour** \\
$E$ & Editors in a shared project, $E \le 4$ \\
$V$ & People with a shared project open \\
$w_h$ & Writes per student per hour** \\
$w_p$ & Writes to one shared project in the period \\
$c$ & Function calls in a month** \\
$e$ & Function outbound data in GB per month \\
$b$ & Hosting downloads in GB per day \\
$h$ & Hosting storage in GB \\
\multicolumn{2}{l}{**Indicates user behaviour} \\
\bottomrule
\end{tabular}
\end{center}

\subsection{Firestore cost per day}

\[
C_{\text{day}} =
\frac{0.09\,\pos{W - 20000}
    + 0.03\,\pos{R - 50000}
    + 0.01\,\pos{D - 20000}}{10^{5}}
\]

Free allowances reset daily and do not roll over, so monthly cost is the sum
of the daily costs:

\[
C_{\text{Firestore}} = \sum_{d=1}^{\text{days}} C_{\text{day}}(d)
\]

\subsection{Activity estimates}

\[
w_h = \underbrace{720}_{\text{presence}}
    + \underbrace{1500}_{\text{typing, estimate}}
    + \underbrace{300}_{\text{saves, estimate}}
    \approx 2500
\]

\[
D \approx 1500\,N H \quad \text{(typing updates removed by cleanup)}
\]

Reads in one shared gallery:

\[
R_{\text{gallery}} = N\,S\,(N - 1)\,H
\]

Reads in one project edited together:

\[
R_{\text{project}} \approx w_p\,(V - 1)
\]

\subsection{Class cost before the free tier}

With $w_h = 2500$ and $S = 300$:

\[
C_{\text{class}} = N H \bigl(0.00225 + 0.00009\,(N - 1)\bigr)
\]

\subsection{Other Firebase costs}

\begin{align*}
C_{\text{storage}} &= 0.15\,\pos{G - 1} && \text{per month} \\
C_{\text{calls}} &= \frac{0.40\,\pos{c - 2000000}}{10^{6}} && \text{per month} \\
C_{\text{outbound}} &= 0.12\,\pos{e - 5} && \text{per month} \\
C_{\text{hosting}} &= 0.15\,\pos{b - 0.36} && \text{per day} \\
C_{\text{hosting storage}} &= 0.026\,\pos{h - 10} && \text{per month}
\end{align*}

Not included: function CPU and memory past 200{,}000 CPU-seconds and
400{,}000 GB-seconds a month, and Authentication past 50{,}000 monthly active
users. Those rates were not checked.

\section{Fixed Costs}
\subsection{Cap Costs}
Many aspects within Wordplay operate within the free tier, fixed costs above the free tier is as follows. 

\subsection{What counts as fixed}

These are paid whether or not anyone uses the app. Most scale with the number
of Firebase projects (\texttt{wordplay-prod} and \texttt{wordplay-dev}) or with
the total size of the database, not with activity. Claude has no fixed fee.

\subsection{Variables}

\begin{center}
\begin{tabular}{ll}
\toprule
Symbol & Meaning \\
\midrule
$P$ & Firebase projects with the functions deployed \\
$J$ & Scheduled jobs on the billing account \\
$j$ & Scheduled jobs per project ($j = 8$) \\
$s$ & Active secret versions on the billing account \\
$a$ & Secret accesses in a month \\
$G$ & Firestore database size in GiB \\
$B_d$ & Daily backups kept ($B_d = 7$ in prod) \\
$B_w$ & Weekly backups kept ($B_w = 14$ in prod) \\
$G_r$ & Function container images stored, in GiB \\
$m$ & Cloud Build minutes in a day \\
\bottomrule
\end{tabular}
\end{center}

\subsection{Scheduled jobs (Cloud Scheduler)}

Each job costs \$0.10 per 31 days. The first 3 jobs on the billing account are
free, and executions are not billed.

\[
J = j\,P, \qquad
C_{\text{scheduler}} = 0.10\,\pos{J - 3} \quad \text{per month}
\]

With $j = 8$ and $P = 2$: $C_{\text{scheduler}} = 0.10 \times 13 = \$1.30$.

Running the jobs costs about 1{,}000 function calls per project each month
(mostly the hourly \texttt{compactProjectUpdates} and the 4-hourly chat
digest), well inside the free 2 million.

\subsection{Secrets (Secret Manager)}

\[
C_{\text{secrets}} = 0.06\,\pos{s - 6}
                   + \frac{0.03\,\pos{a - 10000}}{10^{4}}
                   \quad \text{per month}
\]

Three secrets in two projects is $s = 6$, which is free. A secret is read when
a function instance starts, so $a$ grows with cold starts.

\subsection{Backups and point-in-time recovery (prod only)}

Assuming each backup is the size of the database:

\[
C_{\text{backups}} = 0.03\,(B_d + B_w)\,G = 0.63\,G
\qquad
C_{\text{PITR}} = 0.15\,G
\quad \text{per month}
\]

Protecting the data costs about five times what storing it does. Ignoring the
free 1\,GiB, the total for data is:

\[
C_{\text{data}} \approx 0.15\,G + 0.63\,G + 0.15\,G = 0.93\,G
\quad \text{per month}
\]

These have no free tier and grow with total data kept, not with activity.

\subsection{Deploys}

Every push to \texttt{main} redeploys hosting, rules and every function.

\begin{align*}
C_{\text{build}} &= 0.003\,\pos{m - 120} && \text{per day} \\
C_{\text{registry}} &= 0.10\,\pos{G_r - 0.5} && \text{per month}
\end{align*}

$G_r$ depends on whether old images are cleaned up, which is set in Artifact
Registry rather than in the repository. Hosting files are about 223\,MB per
release, inside the free 10\,GB unless many old releases are kept.

\subsection{Outside Google Cloud}

\begin{center}
\begin{tabular}{ll}
\toprule
Item & Cost \\
\midrule
Resend plan & Not checked; free tier or a monthly plan \\
\texttt{wordplay.dev} domain & Yearly, from the registrar \\
GitHub & Free \\
Anthropic & No fixed fee; pay per token only \\
\bottomrule
\end{tabular}
\end{center}

\subsection{Total fixed cost}

\[
C_{\text{fixed}} =
C_{\text{scheduler}} + C_{\text{secrets}} + C_{\text{backups}}
+ C_{\text{PITR}} + C_{\text{registry}} + \sum_{\text{days}} C_{\text{build}}
+ C_{\text{Resend}} + C_{\text{domain}}
\]

For a database of a few GiB this is a few dollars a month, plus Resend and the
domain.

\section{Other costs}

\subsection{Variables}

\begin{center}
\begin{tabular}{ll}
\toprule
Symbol & Meaning \\
\midrule
$A$ & reCAPTCHA assessments in a month, across the organization \\
$d$ & Active devices that call a Cloud Function in a month \\
$e_w$ & Outbound GB from \texttt{getWebpage} in a month \\
$L$ & New languages translated in a month \\
$G_{\text{restore}}$ & GiB restored from a backup \\
$r_m$ & Documents read by a migration \\
$w_m$ & Documents written by a migration \\
\bottomrule
\end{tabular}
\end{center}

\subsection{App Check (reCAPTCHA)}

A token lasts 7 days and is only requested when a Cloud Function is called,
so each active device makes about one assessment a week:

\[
A \approx 4.3\,d
\]

\[
C_{\text{App Check}} =
\begin{cases}
0 & A \le 10000 \\
8 & 10000 < A \le 100000 \\
8 + \dfrac{A - 100000}{1000} & A > 100000
\end{cases}
\quad \text{per month}
\]

The pricing page lists \$1 per 1{,}000 above 100{,}000 assessments without
saying whether the \$8 fee still applies; this formula assumes it does.

\subsection{\texttt{getWebpage}}

The function fetches any URL with no sign-in, App Check or size limit, so its
cost is set by whoever calls it. Its outbound data shares the function
allowance of 5\,GB, so once that is used each extra GB costs:

\[
C_{\text{getWebpage}} \approx 0.12\,e_w
\quad \text{per month, plus running time}
\]

\subsection{Development use of Claude}

These use the project's Anthropic key from developers' machines and GitHub
Actions, and follow how much the team works rather than how many users there
are:

\[
C_{\text{dev}} \approx 3\,L + C_{\text{changelog}} + C_{\text{scripts}}
\quad \text{per month}
\]

$C_{\text{changelog}}$ is 29 translations per new changelog entry, and
$C_{\text{scripts}}$ covers occasional Opus scripts such as
\texttt{condense-changelog}. Neither was measured.

\subsection{One-off costs}

\begin{align*}
C_{\text{restore}} &= 0.20\,G_{\text{restore}} \\
C_{\text{migration}} &= \frac{0.03\,r_m + 0.09\,w_m}{10^{5}}
\end{align*}

A migration that reads and rewrites 100{,}000 documents costs about \$0.12.
Testing on \texttt{test.wordplay.dev} also uses real Firestore, function and
email costs, at the rates in Section~2.

\subsection{Not modeled}

\begin{center}
\begin{tabular}{ll}
\toprule
Item & Note \\
\midrule
Cloud Logging & Free allowance per project; price above it not checked \\
Firestore outbound data & Free up to 10\,GiB a month; price above it not checked \\
Security rule reads & One extra read when a rule checks a gallery document \\
Crawlers & Run \texttt{getPagePreview} and \texttt{getSitemap} like visitors \\
GitHub Actions & Free while the repository is public \\
Email inbox for \texttt{hi@} & Depends on the provider \\
Tax & Depends on billing location \\
\bottomrule
\end{tabular}
\end{center}

\subsection{Total}

\[
C_{\text{total}} =
C_{\text{Claude}} + C_{\text{Firebase}} + C_{\text{fixed}}
+ C_{\text{App Check}} + C_{\text{getWebpage}} + C_{\text{dev}}
+ C_{\text{one-off}}
\]

\end{document}

